\chapter{Θεωρητικό Υπόβαθρο και Επισκόπηση Βιβλιογραφίας}
\label{chap:background}

Το κεφάλαιο αυτό συγκεντρώνει το μαθηματικό και αλγοριθμικό υπόβαθρο στο οποίο στηρίζεται η παρούσα εργασία, μαζί με την επισκόπηση της σχετικής βιβλιογραφίας. Πρώτα περιγράφονται η δομή του κύβου και η ομάδα των καταστάσεών του, η σημειογραφία των κινήσεων και η μετρική κόστους, καθώς και οι συνθήκες επιλυσιμότητας. Έπειτα παρουσιάζεται το πλαίσιο της ευρετικής αναζήτησης, από τον αλγόριθμο \eng{A*} έως τους πίνακες προτύπων. Η βιβλιογραφική επισκόπηση τοποθετεί στη συνέχεια τους αλγορίθμους \eng{Thistlethwaite}, \eng{Kociemba} και \eng{Korf}, τον αριθμό του Θεού (\eng{God's number}) και τη γενεαλογία των σύγχρονων βέλτιστων επιλυτών. Το κεφάλαιο κλείνει με την τοποθέτηση της παρούσας εργασίας σε αυτό το πλαίσιο.

\section{Ο κύβος και η ομάδα καταστάσεων}
\label{sec:background-group}

Ο κύβος 3×3×3 αποτελείται από 26 ορατά τεμάχια: οκτώ γωνιακά κυβίδια, δώδεκα ακμικά κυβίδια και έξι σταθερά κέντρα. Κάθε στροφή έδρας μεταθέτει τα κυβίδια της έδρας και, ανάλογα με την έδρα, μεταβάλλει τον προσανατολισμό τους. Τα κέντρα δεν μετακινούνται, επομένως ορίζουν σταθερό σύστημα αναφοράς για τα χρώματα.

Για την περιγραφή μιας κατάστασης υπάρχουν δύο φυσικά επίπεδα. Η περιγραφή σε επίπεδο ψηφίδων (\eng{facelets}) καταγράφει τα χρώματα σε 54 θέσεις, από τις οποίες οι 48 μετακινούνται, αφού οι έξι κεντρικές παραμένουν σταθερές. Η περιγραφή σε επίπεδο κυβιδίων προσδιορίζει για κάθε γωνία και κάθε ακμή ποιο τεμάχιο είναι, πού βρίσκεται και πώς είναι προσανατολισμένο: οι οκτώ γωνίες φέρουν μετάθεση και προσανατολισμό modulo 3, ενώ οι δώδεκα ακμές μετάθεση και προσανατολισμό modulo 2. Η πρώτη μορφή είναι πιο κοντά στο πώς αντιλαμβάνεται τον κύβο ο άνθρωπος και κατάλληλη για είσοδο και έξοδο. Η δεύτερη είναι καταλληλότερη για αλγορίθμους, επειδή κάθε κίνηση εκφράζεται ως μετάθεση με μικρή μεταβολή προσανατολισμών, αντί για αλλαγή χρωμάτων στις 48 κινητές θέσεις. Η συγκεκριμένη κωδικοποίηση που υιοθετεί η υλοποίηση περιγράφεται στο Κεφάλαιο~\ref{chap:model}.

Μαθηματικά, οι επιτρεπτές κινήσεις παράγουν πεπερασμένη ομάδα μεταθέσεων με πρόσθετους περιορισμούς προσανατολισμού. Η θεώρηση αυτή είναι κλασική στη βιβλιογραφία του κύβου και στα υπολογιστικά συστήματα άλγεβρας \cite{joyner2008adventures,sageCubeGroup,gapRubikExample}. Η παρούσα εργασία δεν χρησιμοποιεί εξωτερικό σύστημα άλγεβρας για τα αποτελέσματά της. Οι αναφορές αυτές δείχνουν ότι η μοντελοποίηση με μεταθέσεις, υποομάδες και περιορισμούς δεν είναι αυθαίρετη επιλογή του κώδικα, αλλά καθιερωμένος τρόπος περιγραφής του προβλήματος.

Το μέγεθος της ομάδας καταστάσεων \(G\) προκύπτει με ένα απλό συνδυαστικό επιχείρημα. Αν ο κύβος αποσυναρμολογηθεί και συναρμολογηθεί ελεύθερα, οι γωνίες τοποθετούνται με \(8!\) τρόπους και προσανατολίζονται με \(3^{8}\), ενώ οι ακμές τοποθετούνται με \(12!\) τρόπους και προσανατολίζονται με \(2^{12}\). Από τις συναρμολογήσεις αυτές, μόνο το ένα δωδέκατο μπορεί να προκύψει με επιτρεπτές κινήσεις: όπως εξηγείται στην Ενότητα~\ref{sec:background-solvability}, οι περιορισμοί προσανατολισμού γωνιών, προσανατολισμού ακμών και ισοτιμίας εισάγουν παράγοντες 3, 2 και 2 αντίστοιχα. Επομένως
\[
|G| \;=\; \frac{8!\cdot 3^{8}\cdot 12!\cdot 2^{12}}{12}
    \;=\; \frac{8!\cdot 12!\cdot 3^{7}\cdot 2^{11}}{2}
    \;=\; 43\,252\,003\,274\,489\,856\,000 \;\approx\; 4.3\times 10^{19}
\]
νόμιμες καταστάσεις \cite{joyner2008adventures}. Ο αριθμός αυτός είναι το βασικό κίνητρο του θέματος: αποκλείει την απλοϊκή εξαντλητική απαρίθμηση ως πρακτική στρατηγική και καθιστά αναγκαία τη δομημένη αναζήτηση με ισχυρά κάτω φράγματα.

\section{Σημειογραφία κινήσεων και μετρικές}
\label{sec:background-notation}

Η σημειογραφία των κινήσεων ακολουθεί την καθιερωμένη σύμβαση \eng{Singmaster} \cite{singmaster1981notes}. Κάθε έδρα συμβολίζεται με το αρχικό της αγγλικής ονομασίας της: \(U\) (\eng{Up}) είναι η άνω έδρα, \(D\) (\eng{Down}) η κάτω, \(L\) (\eng{Left}) η αριστερή, \(R\) (\eng{Right}) η δεξιά, \(F\) (\eng{Front}) η εμπρόσθια και \(B\) (\eng{Back}) η οπίσθια. Σκέτο γράμμα σημαίνει δεξιόστροφη στροφή της έδρας κατά ένα τέταρτο (90 μοίρες), όπως τη βλέπει κάποιος που την κοιτάζει κατά μέτωπο. Ο τόνος δηλώνει αριστερόστροφη στροφή κατά ένα τέταρτο (π.χ. \(U'\)), ενώ η κατάληξη 2 δηλώνει μισή στροφή 180 μοιρών (π.χ. \(U2\)). Προκύπτουν έτσι 18 σύμβολα κινήσεων, τρία για καθεμία από τις έξι έδρες. Το Σχήμα~\ref{fig:cube-notation} συνοψίζει το ανάπτυγμα του κύβου και τη σημειογραφία.

\begin{figure}[htbp]\centering
% Ανάπτυγμα του κύβου 3×3×3 (διάταξη σταυρού) με τη σημειογραφία κινήσεων.
% Περιεχόμενο σχήματος μόνο: το περιβάλλον figure/caption ορίζεται στο κεφάλαιο.
\begin{tikzpicture}[x=0.55cm, y=0.55cm, line cap=round, line join=round]
  % Οι έξι έδρες του αναπτύγματος: U πάνω, L-F-R-B στη μεσαία ζώνη, D κάτω.
  \foreach \ox/\oy/\lab/\gr in {%
      3/6/U/{πάνω},
      0/3/L/{αριστερά},
      3/3/F/{μπροστά},
      6/3/R/{δεξιά},
      9/3/B/{πίσω},
      3/0/D/{κάτω}}{
    \fill[gray!10] (\ox+1,\oy+1) rectangle (\ox+2,\oy+2);
    \draw[step=1, gray!55, very thin] (\ox,\oy) grid (\ox+3,\oy+3);
    \draw[black, line width=0.8pt] (\ox,\oy) rectangle (\ox+3,\oy+3);
    \node[font=\large\bfseries] at (\ox+1.5,\oy+1.5) {\lab};
    \node[font=\scriptsize, text=black!75, fill=white, inner sep=1pt]
      at (\ox+1.5,\oy+0.5) {\gr};
  }
  % Υπόμνημα: τα τρία επιθήματα κινήσεων της σημειογραφίας Singmaster.
  \node[anchor=west, font=\footnotesize\bfseries] at (13.2,6.6)
    {Επιθήματα κινήσεων};
  \node[anchor=west, font=\footnotesize] at (13.2,5.8)
    {όπου \texttt{X} \(\in\) \{\texttt{U}, \texttt{L}, \texttt{F}, \texttt{R}, \texttt{B}, \texttt{D}\}};
  \node[anchor=west, font=\footnotesize] at (13.2,4.7)
    {\texttt{X}: δεξιόστροφη στροφή \(90^{\circ}\)};
  \node[anchor=west, font=\footnotesize] at (13.2,3.85)
    {\texttt{X\textquotesingle}: αριστερόστροφη στροφή \(90^{\circ}\)};
  \node[anchor=west, font=\footnotesize] at (13.2,3.0)
    {\texttt{X2}: στροφή \(180^{\circ}\)};
\end{tikzpicture}

\caption{Το ανάπτυγμα του κύβου 3×3×3 με τις έξι έδρες και τη σημειογραφία κινήσεων κατά Singmaster: U, D, L, R, F, B για δεξιόστροφες στροφές κατά ένα τέταρτο, με τόνο (π.χ. U\textquotesingle) για αριστερόστροφες και με κατάληξη 2 (π.χ. U2) για μισές στροφές.}
\label{fig:cube-notation}
\end{figure}

Το κόστος μιας ακολουθίας κινήσεων εξαρτάται από τη μετρική. Η παρούσα εργασία χρησιμοποιεί παντού τη μετρική μισών στροφών (\eng{half-turn metric}, HTM), γνωστή και ως \eng{face-turn metric}, στην οποία και οι 18 κινήσεις έχουν κόστος ένα. Η επιλογή αυτή έχει σημασία: στη μετρική τετάρτων στροφής (\eng{quarter-turn metric}, QTM) μια κίνηση όπως η \(R2\) μετρά ως δύο τέταρτα στροφής, οπότε η ίδια ακολουθία αποκτά διαφορετικό μήκος. Κάθε ισχυρισμός για μήκος λύσης ή για απόσταση πρέπει επομένως να δηλώνει τη μετρική του. Η χρήση της HTM συμφωνεί με τη διατύπωση του θέματος και με τη βιβλιογραφία του αριθμού του Θεού \cite{rokicki2010godsnumber,rokicki2013diameter}.

Η μετρική επηρεάζει και τον τρόπο που δομείται η αναζήτηση. Η αναζήτηση κατά πλάτος (\eng{breadth-first search}, BFS) σε βάθος \(d\) εξετάζει όλες τις ακολουθίες μήκους έως \(d\) στο δηλωμένο σύνολο κινήσεων, ενώ η \eng{IDA*} αυξάνει τα όρια βάθους με το ίδιο κόστος. Αλλαγή μετρικής θα σήμαινε διαφορετικό σύνολο βασικών κινήσεων, διαφορετικούς πίνακες και μη συγκρίσιμα αποτελέσματα. Για τον λόγο αυτό η μετρική αποτελεί μέρος του ορισμού κάθε πειράματος της εργασίας.

\section{Επιλυσιμότητα}
\label{sec:background-solvability}

Δεν αντιστοιχεί κάθε τοποθέτηση κυβιδίων σε κατάσταση που μπορεί να προκύψει από επιτρεπτές κινήσεις. Μια διάταξη είναι επιλύσιμη αν και μόνο αν ικανοποιεί τέσσερις συνθήκες \cite{joyner2008adventures}. Κάθε γωνία και κάθε ακμή εμφανίζεται ακριβώς μία φορά. Το άθροισμα των προσανατολισμών των γωνιών είναι μηδέν modulo 3. Το άθροισμα των αναστροφών (\eng{flips}) των ακμών είναι μηδέν modulo 2. Η ισοτιμία της μετάθεσης των γωνιών συμπίπτει με την ισοτιμία της μετάθεσης των ακμών. Η αναγκαιότητα των συνθηκών προκύπτει επειδή η λυμένη κατάσταση τις ικανοποιεί και κάθε βασική κίνηση τις διατηρεί. Η επάρκειά τους αποδεικνύεται κατασκευαστικά στη βιβλιογραφία της θεωρίας ομάδων του κύβου \cite{joyner2008adventures}.

Οι συνθήκες αυτές έχουν άμεσες πρακτικές συνέπειες. Μία μόνο στραμμένη γωνία ή μία μόνο αναστραμμένη ακμή παραβιάζει τους περιορισμούς προσανατολισμού, ενώ ασυμφωνία ισοτιμίας σημαίνει ότι η κατάσταση δεν παράγεται από νόμιμες κινήσεις. Έτσι, μια περιγραφή με σωστά πλήθη χρωμάτων μπορεί παρ' όλα αυτά να μην αντιστοιχεί σε φυσικά επιλύσιμη κατάσταση. Άρα ο έλεγχος εγκυρότητας πρέπει να αφορά τη φυσική επιλυσιμότητα, δηλαδή τη δομή του κύβου, και όχι μόνο τη συντακτική ορθότητα και το πλήθος των χρωμάτων. Ένας επιλυτής δεν πρέπει να σπαταλά χρόνο σε καταστάσεις εκτός του χώρου του κύβου· γι' αυτό κάθε πειραματική είσοδος ελέγχεται ως προς τη νομιμότητά της. Ο τρόπος με τον οποίο η υλοποίηση πραγματοποιεί τους ελέγχους αυτούς περιγράφεται στο Κεφάλαιο~\ref{chap:model}.

Οι τρεις τελευταίες συνθήκες εξηγούν και τον παράγοντα 12 της Ενότητας~\ref{sec:background-group}: οι ελεύθερες συναρμολογήσεις χωρίζονται σε \(3\cdot2\cdot2=12\) κλάσεις ίσου μεγέθους, από τις οποίες μόνο μία αποτελεί τον χώρο των νόμιμων καταστάσεων.

\section{Αναζήτηση και ευρετικές συναρτήσεις}
\label{sec:background-search}

Η εύρεση ελάχιστης λύσης είναι πρόβλημα αναζήτησης συντομότερου μονοπατιού σε έναν τεράστιο γράφο που δεν τον έχουμε αποθηκευμένο ολόκληρο. Η ορολογία του χώρου καταστάσεων (\eng{state space}), του μετώπου αναζήτησης (\eng{search frontier}) και του ασυμπτωτικού κόστους ακολουθεί τη γενική βιβλιογραφία αλγορίθμων και τεχνητής νοημοσύνης \cite{russellNorvig2020aima,cormen2022clrs}.

Η BFS έχει μια ιδιότητα καθοριστική για τη βελτιστότητα. Όταν ολοκληρώνει όλα τα επίπεδα μέχρι ένα βάθος και βρίσκει την πρώτη λύση στο επίπεδο \(d\), αποδεικνύει άμεσα ότι δεν υπάρχει συντομότερη λύση, αφού όλα τα μικρότερα βάθη έχουν ήδη εξεταστεί. Το μειονέκτημα είναι η μνήμη: το μέτωπο αναζήτησης και το σύνολο των επισκεμμένων καταστάσεων αυξάνονται εκθετικά με το βάθος. Στον κύβο 3×3×3 η BFS είναι πρακτική μόνο για ρηχά βάθη και για ελέγχους σύγκρισης.

Ο αλγόριθμος A* των \eng{Hart}, \eng{Nilsson} και \eng{Raphael} θεμελίωσε το πλαίσιο της ευρετικής αναζήτησης ελάχιστου κόστους \cite{hart1968astar}. Κεντρική έννοια είναι η παραδεκτή ευρετική (\eng{admissible heuristic}): μια συνάρτηση \(h\) που δεν υπερεκτιμά ποτέ το πραγματικό υπόλοιπο κόστος. Με παραδεκτή ευρετική, η λύση που επιστρέφει ο A* είναι αποδεδειγμένα βέλτιστη. Η συστηματική ανάλυση των ευρετικών συναρτήσεων αποτελεί κεντρικό αντικείμενο της κλασικής βιβλιογραφίας \cite{pearl1984heuristics}. Ο A* όμως διατηρεί στη μνήμη το σύνολο των ανοιχτών κόμβων, κάτι ανέφικτο για χώρους της τάξης του \(10^{19}\).

Η IDA* του Korf συνδυάζει την επαναληπτική εμβάθυνση (\eng{iterative deepening}) με την παραδεκτή ευρετική και απαιτεί μνήμη γραμμική ως προς το βάθος \cite{korf1985iddfs}. Η αναζήτηση εκτελεί διαδοχικές διελεύσεις βάθους με όριο στην τιμή \(f(n)=g(n)+h(n)\), όπου \(g(n)\) το κόστος έως τον τρέχοντα κόμβο και \(h(n)\) η ευρετική εκτίμηση του υπολοίπου. Η βελτιστότητα μπορεί να διατυπωθεί με ακρίβεια ως εξής: αν η \(h\) είναι παραδεκτή και το κατώφλι αυξάνεται κάθε φορά στην ελάχιστη τιμή \(f\) που υπερέβη το προηγούμενο όριο, τότε η πρώτη λύση που εντοπίζεται είναι αποδεδειγμένα ελάχιστη. Η πρακτική απόδοση όμως εξαρτάται καθοριστικά από την ποιότητα της ευρετικής. Με αδύναμο κάτω φράγμα, η IDA* εκφυλίζεται σε απλή επαναληπτική εμβάθυνση και εξετάζει υπέρμετρα μεγάλο πλήθος κόμβων.

Το απλούστερο παραδεκτό κάτω φράγμα για τον κύβο προκύπτει αν διαιρέσουμε με το τέσσερα τον αριθμό των κυβιδίων που βρίσκονται σε λάθος θέση ή προσανατολισμό και στρογγυλοποιήσουμε προς τα πάνω. Μία κίνηση μετακινεί το πολύ τέσσερις γωνίες και τέσσερις ακμές και αλλάζει τον προσανατολισμό το πολύ τεσσάρων γωνιών ή τεσσάρων ακμών. Επομένως κάθε τέτοιο πηλίκο, δηλαδή το πλήθος δια τέσσερα, είναι παραδεκτό κάτω φράγμα, όπως και το μέγιστό τους. Το φράγμα αυτό όμως είναι συνήθως πολύ χαμηλό για βαθιές καταστάσεις.

Ισχυρότερα φράγματα δίνουν οι πίνακες προτύπων (\eng{pattern databases}, PDB), τους οποίους εισήγαγαν οι \eng{Culberson} και \eng{Schaeffer} \cite{culberson1996pdb}. Αποθηκεύουν τις ακριβείς αποστάσεις ενός απλοποιημένου υποπροβλήματος και χρησιμοποιούνται ως παραδεκτή ευρετική για το πλήρες πρόβλημα. Στον κύβο, η αφαίρεση εκφράζεται ως προβολή σε συντεταγμένη (\eng{coordinate}): μια ακέραια κωδικοποίηση μιας πιο περιορισμένης πτυχής της κατάστασης, όπως ο προσανατολισμός των γωνιών, ο προσανατολισμός των ακμών ή οι θέσεις των τεσσάρων ακμών της μεσαίας φέτας (\eng{UD-slice}). Κάθε ακολουθία που λύνει την πλήρη κατάσταση λύνει κατ' ανάγκη και καθεμία από τις προβολές της. Η απόσταση μιας προβολής από τη λυμένη μορφή της είναι επομένως παραδεκτό κάτω φράγμα της πλήρους απόστασης, και το μέγιστο πολλών τέτοιων φραγμάτων παραμένει παραδεκτό. Χαρακτηριστικές κλίμακες είναι η συνδυασμένη γωνιακή προβολή με \(8!\cdot3^{7}=88\,179\,840\) καταστάσεις και καθεμία από τις προβολές έξι ακμών με \(42\,577\,920\) καταστάσεις. Οι αποστάσεις κάθε προβολής υπολογίζονται πλήρως με BFS στον αντίστοιχο μικρότερο γράφο και αποθηκεύονται είτε ως πίνακες προτύπων είτε ως πίνακες αποκοπής (\eng{pruning tables}) που καθοδηγούν την αποκοπή κλάδων της αναζήτησης. Οι συγκεκριμένες κωδικοποιήσεις της εργασίας περιγράφονται στο Κεφάλαιο~\ref{chap:model}.

Οι \eng{Felner}, Korf και \eng{Hanan} μελέτησαν αργότερα τους προσθετικούς πίνακες προτύπων (\eng{additive PDB}) \cite{felner2004additivePDB}. Όταν οι αφαιρέσεις αφορούν ξένα μεταξύ τους σύνολα κυβιδίων και κάθε πίνακας μετρά μόνο τις κινήσεις των δικών του τεμαχίων, οι τιμές μπορούν να αθροιστούν αντί να μεγιστοποιηθούν, δίνοντας ισχυρότερο παραδεκτό φράγμα. Η διάκριση ανάμεσα στον συνδυασμό με μέγιστο και στον προσθετικό συνδυασμό επανέρχεται στην Ενότητα~\ref{sec:background-literature}.

Χρειάζεται προσοχή στην ερμηνεία των φραγμάτων. Αν ένας πίνακας δίνει κάτω φράγμα 5, αυτό σημαίνει ότι τουλάχιστον μία προβολή απαιτεί πέντε κινήσεις. Δεν σημαίνει ότι η πλήρης απόσταση είναι 5, επειδή οι προβολές αγνοούν πληροφορία. Γενικότερα, το τι μπορούμε να ισχυριστούμε εξαρτάται από το ποια διαδικασία ολοκληρώθηκε και πόσο πλήρως. Ολοκληρωμένη BFS μέχρι το βάθος της πρώτης λύσης αποδεικνύει ακριβή απόσταση. Ολοκληρωμένη IDA* με παραδεκτή ευρετική αποδεικνύει επίσης ελάχιστο μήκος για τη συγκεκριμένη κατάσταση. Ένας σταδιακός επιλυτής που επιστρέφει επαληθευμένη λύση τεκμηριώνει την ύπαρξη λύσης, όχι την ελαχιστότητά της. Μια αναζήτηση που διακόπηκε πριν ολοκληρωθεί τεκμηριώνει μόνο κάτω φράγμα ή την απουσία απόδειξης εντός των ορίων. Η ακριβής ορολογία με την οποία καταγράφονται αυτές οι διαβαθμίσεις στα αποτελέσματα ορίζεται στο Κεφάλαιο~\ref{chap:validation}.

\section{Επισκόπηση βιβλιογραφίας}
\label{sec:background-literature}

Η επισκόπηση δεν επιχειρεί να καλύψει όλους τους γνωστούς τρόπους επίλυσης του κύβου. Ο στόχος της είναι στενότερος: να τοποθετήσει τους αλγορίθμους Thistlethwaite, Kociemba και Korf στο πλαίσιο της αναζήτησης και των παραδεκτών ευρετικών, να τεκμηριώσει τον αριθμό του Θεού ως εξωτερικό ορόσημο και να παρουσιάσει τη γενεαλογία των σύγχρονων βέλτιστων επιλυτών. Σε αυτήν εντάσσεται το ισχυρότερο τεκμηριωτικό σκέλος της εργασίας.

\subsection{Thistlethwaite και η αλυσίδα υποομάδων}
\label{sec:background-thistlethwaite}

Ο αλγόριθμος Thistlethwaite εισήγαγε τη μεθοδική σταδιακή μείωση του χώρου αναζήτησης μέσω αλυσίδας υποομάδων. Αντί να αναζητείται άμεση λύση από οποιαδήποτε κατάσταση, επιβάλλονται διαδοχικοί περιορισμοί που οδηγούν την κατάσταση σε ολοένα μικρότερες υποομάδες. Σε κάθε στάδιο επιτρέπονται μόνο κινήσεις που διατηρούν τους περιορισμούς των προηγούμενων σταδίων, μέχρι ο εναπομένων χώρος να γίνει αρκετά μικρός ώστε να ολοκληρωθεί η λύση. Η τεχνική περιγραφή του \eng{Scherphuis} και η τεκμηρίωση του Kociemba παρουσιάζουν τη μέθοδο ως αλυσίδα φάσεων με πίνακες κινήσεων και γνωστό άνω φράγμα συνολικών κινήσεων \cite{scherphuisThistlethwaite,kociembaImplementationDetails}.

Η μέθοδος έχει κυρίως εννοιολογική αξία: κάνει ορατή την έννοια της υποομάδας και δείχνει ότι η σχεδίαση του χώρου αναζήτησης μπορεί να αντικαταστήσει την καθοδήγηση μέσω ευρετικής. Η παρούσα εργασία υλοποιεί σε C μια τετραφασική εκδοχή της αλυσίδας, καθοδηγούμενη από τέσσερις ακριβείς πίνακες αποστάσεων συμπλόκων (\eng{cosets}) που παράγονται με BFS. Οι τέσσερις φάσεις είναι: προσανατολισμός ακμών· μηδενισμός προσανατολισμού γωνιών μαζί με τοποθέτηση των ακμών της μεσαίας φέτας στη φέτα τους· είσοδος στην υποομάδα \(G_3\) των μισών στροφών· τελική λύση μόνο με μισές στροφές. Η υλοποίηση περιγράφεται στο Κεφάλαιο~\ref{chap:algorithms}. Οι λύσεις της δεν εγγυώνται ελαχιστότητα, επειδή δεν υπάρχει απόδειξη ότι η συνολική ακολουθία είναι η συντομότερη δυνατή.

\subsection{Kociemba και ο αλγόριθμος δύο φάσεων}
\label{sec:background-kociemba}

Ο αλγόριθμος δύο φάσεων του Kociemba συμπτύσσει στην πράξη τη σταδιακή μείωση σε δύο μόνο βήματα \cite{kociembaTwoPhase}. Η πρώτη φάση φέρνει την κατάσταση στην ενδιάμεση υποομάδα όπου οι προσανατολισμοί έχουν μηδενιστεί και οι τέσσερις ακμές της μεσαίας φέτας βρίσκονται στη φέτα τους. Η δεύτερη φάση ολοκληρώνει τη λύση μέσα σε περιορισμένο σύνολο κινήσεων. Η τεκμηρίωση του Kociemba περιγράφει τις συντεταγμένες, τους πίνακες κινήσεων και τους πίνακες αποκοπής που καθιστούν τη μέθοδο εξαιρετικά αποδοτική \cite{kociembaImplementationDetails}, ενώ το πρόγραμμα Cube Explorer αποτελεί το γνωστότερο παράδειγμα ώριμης υλοποίησης \cite{kociembaCubeExplorer}. Η μέθοδος δίνει πολύ σύντομες λύσεις σε ελάχιστο χρόνο, χωρίς όμως να εγγυάται ότι κάθε λύση που επιστρέφει είναι η συντομότερη δυνατή.

Η παρούσα εργασία διακρίνει τρία επίπεδα: τη βιβλιογραφική μέθοδο, έναν εξωτερικό προσαρμογέα (\eng{adapter}) και μια εγγενή υλοποίηση περιορισμένου εύρους. Ο προσαρμογέας βασίζεται στο πακέτο \code{kociemba} του δημόσιου καταλόγου πακέτων PyPI, γραμμένο από τον \eng{muodov} και διαθέσιμο με άδεια GPLv2 \cite{muodovKociembaPkg}, και χρησιμοποιείται ως πρακτικό σημείο αναφοράς για σύγκριση. Η εγγενής υλοποίηση χρησιμοποιεί πραγματικούς πίνακες προβολών για τη φάση 1 και πρόσθετες φασικές προβολές για τη φάση 2· περιγράφεται στο Κεφάλαιο~\ref{chap:algorithms} και αποτελεί την τοπική αλγοριθμική συνεισφορά. Η εργασία δεν την παρουσιάζει ως ισοδύναμη με το Cube Explorer, επειδή λείπουν ισχυρότεροι συνδυασμένοι πίνακες και εκτενέστερες βελτιστοποιήσεις.

\subsection{Korf, IDA* και οι πίνακες προτύπων στον κύβο}
\label{sec:background-korf}

Ο Korf έδωσε τις πρώτες βέλτιστες λύσεις τυχαίων καταστάσεων του κύβου 3×3×3, συνδυάζοντας την IDA* με πίνακες προτύπων \cite{korf1997pattern}. Η ευρετική του προέκυπτε από έναν πλήρη γωνιακό πίνακα και δύο πίνακες έξι ακμών, παίρνοντας το μέγιστο των τριών τιμών. Ο συνδυασμός με μέγιστο διατηρεί την παραδεκτότητα χωρίς να απαιτεί αφαιρέσεις πάνω σε ξένα μεταξύ τους σύνολα κυβιδίων. Οι προσθετικοί πίνακες προτύπων των Felner, Korf και Hanan αποτελούν μεταγενέστερη γενίκευση \cite{felner2004additivePDB} και δεν πρέπει να συγχέονται με το αρχικό σχήμα του 1997.

Η παρούσα εργασία ακολουθεί ως προεπιλογή το σχήμα του Korf: παράγει πλήρη γωνιακό πίνακα προτύπων με 88\,179\,840 καταστάσεις και οκτώ πλήρεις πίνακες έξι ακμών, και χρησιμοποιεί ως κύρια ευρετική το μέγιστο των διαθέσιμων παραδεκτών φραγμάτων. Το άθροισμα τιμών από πίνακες που μετρούν το πλήρες κόστος δεν χρησιμοποιείται, επειδή δεν θα ήταν παραδεκτό για επικαλυπτόμενες προβολές. Επιπλέον, υπάρχει μια πειραματική υλοποίηση κατάτμησης κόστους (\eng{cost partitioning}) με δύο συμβατές εξακμικές προβολές URF/DLB και κόστη τελεστών 0/1. Καταγράφεται χωριστά στα τεκμήρια και δεν αντικαθιστά την προεπιλεγμένη ευρετική. Ο γενικότερος και ισχυρότερος προσθετικός συνδυασμός πινάκων προτύπων, πέρα από αυτή την περιορισμένη δοκιμή, μένει ως μελλοντική εργασία.

\subsection{Ο αριθμός του Θεού και η διάμετρος της ομάδας}
\label{sec:background-godsnumber}

Ο αριθμός του Θεού είναι η διάμετρος του γράφου καταστάσεων: το μέγιστο, επί όλων των καταστάσεων, της ελάχιστης απόστασης από τη λυμένη κατάσταση. Το 1995 ο Michael Reid απέδειξε ότι η διάταξη \eng{superflip}, στην οποία όλα τα κυβίδια βρίσκονται στη θέση τους αλλά και οι δώδεκα ακμές είναι αναστραμμένες, δεν μπορεί να λυθεί με λιγότερες από 20 κινήσεις στη μετρική HTM \cite{reid1995superflip}. Αφού υπάρχει και γνωστή λύση 20 κινήσεων για το superflip, η απόστασή του είναι ακριβώς 20· έτσι το αποτέλεσμα έδωσε το πρώτο κάτω φράγμα 20 για τη διάμετρο. Το 1997 ο ίδιος παρουσίασε βέλτιστο επιλυτή που χρησιμοποιεί ως παραδεκτό κάτω φράγμα το μέγιστο της απόστασης από την ενδιάμεση υποομάδα της φάσης 1 και στους τρεις άξονες \cite{reid1997optimal}. Το «φράγμα τριών αξόνων» του Reid αξιοποιείται και από την παρούσα εργασία στην απόδειξη της απόστασης του superflip (Κεφάλαιο~\ref{chap:implementation}).

Το 2010 οι Rokicki, Kociemba, Davidson και Dethridge ολοκλήρωσαν την απόδειξη ότι η διάμετρος είναι ακριβώς 20: κάθε κατάσταση λύνεται σε 20 το πολύ κινήσεις στη μετρική HTM \cite{rokicki2010godsnumber}, με την πλήρη δημοσίευση να ακολουθεί \cite{rokicki2013diameter}. Η απόδειξη απαίτησε τεράστιο υπολογιστικό φόρτο, με αξιοποίηση συμμετριών και ταξινόμηση συμπλόκων. Η παρούσα εργασία χρησιμοποιεί το αποτέλεσμα ως εξωτερικό ορόσημο και δεν ισχυρίζεται ότι το αναπαράγει.

\subsection{Σύγχρονοι βέλτιστοι επιλυτές}
\label{sec:background-modern}

Η εξέλιξη των βέλτιστων επιλυτών δεν σταμάτησε στον Korf και στον Reid. Ο βέλτιστος επιλυτής του Kociemba εφαρμόζει το φράγμα της φάσης 1 και στους τρεις άξονες, στο πνεύμα του Reid, μέσα στο Cube Explorer \cite{kociembaOptimalSolver}. Ο Rokicki, με το πρόγραμμα nxopt, κλιμάκωσε το ίδιο σχήμα σε πολύ μεγαλύτερους πίνακες αποκοπής, θυσιάζοντας περισσότερη μνήμη για ισχυρότερο φράγμα \cite{rokickiNxopt}. Πιο πρόσφατα, ο επιλυτής Nissy του Tronto και η οικογένεια πινάκων H48 του \code{nissy-core} (Tronto και Tenuti) προσφέρουν μια σύγχρονη, ανοιχτού κώδικα υλοποίηση βέλτιστης επίλυσης, με συμμετρικά ανηγμένους πίνακες αποκοπής μεγάλης κλίμακας \cite{trontoNissyCoreH48}.

Η παρούσα εργασία εντάσσεται ρητά σε αυτή τη γενεαλογία: τα ισχυρότερα τεκμήρια ακριβούς αναζήτησης σε μεγάλα βάθη τα παράγει μια ενσωματωμένη μηχανή επίλυσης (\eng{backend}) βασισμένη στο \code{nissy-core}, με πίνακα αποκοπής τύπου \code{h48h7}. Τα αποτελέσματα διασταυρώνονται ανεξάρτητα με τον εξωτερικό επιλυτή Nissy (Κεφάλαια~\ref{chap:implementation} και~\ref{chap:validation}).

\subsection{Ο κύβος 2×2×2 ως πλήρης μικρότερη περίπτωση}
\label{sec:background-pocket}

Ο κύβος 2×2×2 (\eng{Pocket Cube}) έχει σημαντικά μικρότερο, αλλά όχι τετριμμένο, χώρο καταστάσεων· το εκπαιδευτικό υλικό του Randelshofer περιγράφει το πρόβλημα και τις κινήσεις του \cite{randelshoferPocketCube}. Με την καθιερωμένη κανονικοποίηση που κρατά σταθερή τη γωνία DBL, ο χώρος έχει \(7!\cdot3^{6}=3\,674\,160\) καταστάσεις: είναι αρκετά μεγάλος ώστε να μην αποτελεί τετριμμένο παράδειγμα, αλλά αρκετά μικρός ώστε να επιτρέπει πλήρη απαρίθμηση με BFS σε συμβατικό υπολογιστή. Στην παρούσα εργασία ο 2×2×2 δεν υποκαθιστά τον 3×3×3. Χρησιμοποιείται ως πλήρης μελέτη περίπτωσης για το είδος της απόδειξης που χρειάζεται ένας βέλτιστος επιλυτής, με πλήρη κατανομή ακριβών αποστάσεων και αντιπροσωπευτικές λύσεις (Κεφάλαιο~\ref{chap:experiments}).

\subsection{Γενικευμένη υπολογιστική δυσκολία}
\label{sec:background-complexity}

Για τις γενικευμένες οικογένειες κύβων \(n\times n\times n\), οι Demaine και συνεργάτες έδωσαν ασυμπτωτικά όρια για τη διάμετρο και αντίστοιχους αλγορίθμους επίλυσης \cite{demaine2011rubik}, ενώ αποδείχθηκε ότι το πρόβλημα εύρεσης βέλτιστης λύσης είναι NP-πλήρες (\eng{NP-complete}) \cite{demaine2018npcomplete}. Ο σταθερός φυσικός κύβος 3×3×3 έχει πεπερασμένο χώρο, οπότε κάθε ερώτημα απόστασης μπορεί θεωρητικά να απαντηθεί με εξαντλητική αναζήτηση και επαρκή μνήμη. Τα αποτελέσματα αυτά δεν μπορούν επομένως να μεταφερθούν ως ισχυρισμός πολυπλοκότητας για τη σταθερή περίπτωση. Χρησιμοποιούνται ως πλαίσιο που εξηγεί γιατί η βέλτιστη επίλυση δεν είναι απλή άσκηση εξαντλητικής αναζήτησης και γιατί η κλίμακα των πινάκων, των συμμετριών και της αναζήτησης παραμένει κεντρικό ζήτημα.

\section{Τοποθέτηση της παρούσας εργασίας}
\label{sec:background-positioning}

Η εργασία είναι συντηρητική ως προς τη χρήση των βιβλιογραφικών αναφορών: οι αναφορές στηρίζουν ιστορικούς, μαθηματικούς και αλγοριθμικούς ισχυρισμούς και δεν υποκαθιστούν την τοπική τεκμηρίωση. Η διάκριση αυτή προστατεύει από τις υπερβολικές προσδοκίες που μπορεί να δημιουργήσει ένας τίτλος με τη λέξη «βέλτιστη». Το ότι η βιβλιογραφία γνωρίζει τη διάμετρο 20 δεν σημαίνει ότι κάθε υλοποίηση μπορεί να πιστοποιεί βέλτιστες λύσεις για αυθαίρετες καταστάσεις. Αν ένας σταδιακός επιλυτής επιστρέψει επαληθευμένη λύση μήκους 21, αυτό δεν αντιφάσκει με τη διάμετρο 20, επειδή ο επιλυτής δεν ισχυρίζεται ελαχιστότητα. Αν μια ακριβής αναζήτηση διακοπεί, το αποτέλεσμα λέει μόνο ότι δεν ολοκληρώθηκε απόδειξη εντός των δηλωμένων ορίων.

Μέσα σε αυτό το πλαίσιο, η συνεισφορά της εργασίας ορίζεται με σαφήνεια: αναπαραγώγιμη εκπαιδευτική και πειραματική υλοποίηση των τριών κλασικών μεθόδων, παραδεκτά κάτω φράγματα από παραγόμενους πίνακες, επαληθευμένο σώμα μετρήσεων, πλήρης εξαντλητική μελέτη του 2×2×2 και τεκμήρια ακριβούς αναζήτησης σε μεγάλα βάθη μέσω της γενεαλογίας H48, με ανεξάρτητη διασταύρωση. Η απαίτηση του θέματος για εκτέλεση σε συμβατικό υπολογιστή καθόρισε τις αποφάσεις μας ως προς την κλίμακα: προτιμήθηκαν δομές που μπορούν να παραχθούν, να αποθηκευτούν και να ελεγχθούν τοπικά.

Με αυτή τη λογική, καμία αναφορά δεν μπαίνει διακοσμητικά. Η θεωρία ομάδων εξηγεί το μοντέλο και η σημειογραφία τις κινήσεις. Η βιβλιογραφία των A*, IDA* και πινάκων προτύπων εξηγεί την ακριβή αναζήτηση και τα κάτω φράγματα, ενώ οι πηγές Thistlethwaite και Kociemba τις σταδιακές μεθόδους. Τέλος, η βιβλιογραφία του αριθμού του Θεού δίνει το καθολικό ορόσημο και η γενεαλογία των σύγχρονων βέλτιστων επιλυτών το σημείο αναφοράς για τα τεκμήρια μεγάλου βάθους. Τα τοπικά αποτελέσματα τεκμηριώνονται χωριστά, με τα αποθηκευμένα αρχεία αποτελεσμάτων και τους ελέγχους που περιγράφονται στα Κεφάλαια~\ref{chap:validation} και~\ref{chap:experiments}.
